\section{Prepare compressed braid factors only when they are visited}
\label{sec:lazy-forest}

Section~\ref{sec:adaptive-forest} stops recognition after a complete
nontrivial-factor proof, but still constructs and summarizes every projected
factor before choosing the first one. Thus the earlier scheduler saves
leaf recognition work while paying all factor-preparation costs. The
maintained recognizer now constructs a cheaper schedule directly from the
input grammar, and prepares factors on demand. The topology theorem and
certificate formats are unchanged.

\subsection{Keep global component validation before scheduling}

The input grammar is validated in full, including unreachable rules. Its
root occurrence counts determine every singleton-generator cut. Its exact
root permutation must describe a knot before any factor can prove a
nontrivial-knot verdict. A link containing a nontrivial component remains a
link; one obstructing factor does not replace this global check.

As before, deleting singleton generators partitions consecutive strand
intervals and expresses the knot as a connected sum of their projected
closures. One nontrivial summand proves the whole knot nontrivial, whereas
an unknot verdict requires every summand to be an unknot. A selected-factor
proof therefore remains the existing version-three certificate with one
source-bound projected grammar. Positive certificates retain every factor
in original strand order, regardless of discovery order. The checker still
reconstructs the selected projection, or all projections for a complete
positive proof, without consulting the producer's schedule.

\subsection{A schedule without constructing the projected grammars}

Let $s$ be the number of supplied grammar rules, $m$ the number of strands,
and $f$ the number of factor intervals. Complete root support implies
$m\le s+1$, so a strand-label-to-factor array is safe after global knot
validation. Cut labels map to no factor. Summing root occurrence counts
inside each interval gives its exact projected word length.

A reverse DAG pass propagates root weights to children, adding both
contributions when a concatenation repeats the same child. Each terminal
contributes its weight with its sign to the exponent of its factor.
Unreachable terminals have weight zero. Consequently these projected
lengths and exponents are exact without expanding any word or building
any factor grammar.

To estimate construction cost, a forward pass assigns each node either
an empty span or the hull of the factor indices touched by its terminal
descendants. A terminal at a cut has empty span. A concatenation takes the
hull of its nonempty child spans. The hull may contain factors that the
node does not actually touch; this only makes the estimate larger.

A terminal contributes a possible new rule to its single factor. A
concatenation contributes only where both children can survive: the
intersection of the two child hulls. If one child is empty in a factor,
projection aliases the surviving child and creates no rule. An empty child
span or a disjoint pair of child spans therefore contributes nothing.

For each nonempty contribution interval $[a,b]$, add one to a difference
array at $a$ and subtract one at $b+1$. Prefix summation gives the possible
rule creations per factor. Its value plus one bounds the number of rules
in the canonical projection. Every created terminal has the corresponding
factor label; every created concatenation has two nonempty projected
children, so that factor belongs to both child hulls. The extra rule is
the empty node. Interning and false-positive hull intersections can only
make the estimate larger than the actual count. All supplied nodes are
considered, including unreachable ones, because the established canonical
projection includes their surviving rules. This differs deliberately from
root-weighted length and exponent.

The schedule keeps the previous three priority classes: at-most-two-strand
or exponent-obstructed factors first, other three-strand factors next, and
wider residual factors last. Within a class it sorts by this rule-count
upper bound, then original index. The old tie-breaker was the actual
projected rule count, which required eager construction. The new estimate
is a scheduling heuristic, not a topological certificate or a guarantee of
optimal discovery order.

\subsection{Lazy preparation and unchanged replay}

Only when a factor is visited does the producer construct its canonical
projection and run the full leaf summary. It checks the resulting length
and exponent against the schedule and checks that the leaf is a knot.
The existing finite, compressed three-braid, elementary-reduction and
exceptional-cube backends then produce the same kinds of leaf proofs.
A negative verdict immediately discards the remaining schedule; unvisited
factor grammars are never constructed. An inconclusive wider factor does
not justify a negative verdict, and a positive result still requires all
factors.

A one-strand projection has a further exact shortcut: every generator is
erased, so its canonical grammar is just the empty rule with root zero.
This agrees byte-for-byte with the old projection, including in the
presence of unreachable source rules. Producer and verifier can construct
it directly after global source validation. Invalid unreachable rules
still fail the earlier validation.

The public option \texttt{use\_lazy\_factors=False}, or CLI flag
\texttt{--eager-factors}, retains eager preparation and its former exact
rule-count ordering. The one-strand shortcut applies to both policies.
No certificate version change is necessary: scheduling choices are not
trusted proof fields, and every published verdict still undergoes fresh
independent replay. Old proofs replay with the current checker and new
proofs replay with the pinned previous checker.

\subsection{Preparation bound and its limits}

Let $q$ be the number of factors visited before completion, and let $s_i,m_i$
be the rule and strand counts of a visited projection. The plan uses
$O(s+m+f\log f)$ structural operations and $O(s+m+f)$ working entries, with
binary root weights charged according to their bit lengths. A single-factor
input bypasses the scheduling pass entirely. Ordinary projection scans the
$s$ supplied rules; a one-strand projection is constant size. Apart from
leaf recognition and binary arithmetic on lengths and weights, preparation now costs
\[
 O\!\left(s+m+f\log f+\sum_{i\text{ visited}}(s+s_i m_i)\right),
\]
where the second term inside the sum bounds the unchanged full leaf
permutation summary. Previously every factor paid its projection and
summary cost before recognition could stop. The new producer retains only
visited projections and their required evidence, rather than materializing
all factors first.

At the checkpoint measured in this section, the global root's dense-permutation pass costs $O(sm)$ word operations and may dominate a large forest. That work is intentionally
included in every complete-call measurement. This optimization therefore
does not claim to remove the root permutation cost or improve the overall
worst-case asymptotic bound merely by reducing factor preparation. If all
factors are needed, $q=f$, and planning adds work. The exceptional-factor
exponential bound and the restricted quasi-polynomial regimes of the prior
sections remain unchanged.

All planning, preparation, leaf solving, proof transport accounting and
replay share the existing public work/time allowances. The estimate affects
only which work is tried first. A limit still returns inconclusive with no
partial certificate; external cancellation, including \texttt{ValueError},
is not converted to an ordinary resource result.

\subsection{Native validation}

Seventy-eight existing compressed-braid tests pass in 20.043 seconds, and
eight final focused tests pass in 0.244 seconds. The new tests compare exact
projected metadata and structural bounds on 107 grammars, including nonconvex support spans, dead
rules and repeated child references. They check that a thirteen-factor
finite-obstruction example makes exactly two projection calls: one in
production and one in mandatory replay, both for the selected factor.
Positive forests retain identical complete certificates and include every
factor. Link inputs are rejected before planning; one-strand shortcuts
preserve exact canonical output; replay works with the planner disabled;
and exact work limits and cancellation retain their prior semantics.

A further native audit compares 160 prior explicit braid cases and 80
seeded forests against the actual maintained implementation at commit
\texttt{586bfc8d0ffc}. There are 110 unknot and 130 nontrivial-knot cases.
Every old and new certificate replays in both implementations; every
positive certificate is identical. Scheduling metadata and rule bounds
are also checked against actual projections. The 4.424-second final audit keeps
all source grammars, work counts and factor orders in
\path{data/lazy-forest-audit.json}. Its bounded-work example completes a
finite negative proof within an allowance at which the old implementation
is inconclusive, with both string-node and cube-generator allowances zero.

The final full maintained suite passes 940 tests with Regina in 208.823 seconds.
The benchmark compares the pinned old implementation, an identical old
control, current defaults, an identical current control and current eager
preparation. Each of ten fixtures has one shuffled warmup and five measured
rounds. Every timed call is the complete public recognizer: source checks,
root summaries, factor discovery, proof construction, transport allowance
accounting and mandatory independent replay. Projection-call counts come
from separate untimed instrumented calls. No precomputed projections are
supplied to a timed arm.

\begin{center}\small
\begin{tabular}{@{}lrrrrr@{}}
\toprule
Input & old ms & lazy ms & eager ms & ratio & projections old/new \\
\midrule
late\_finite\_8 & 14.092 & 11.469 & 14.063 & 1.255 & 10/2 \\
late\_finite\_24 & 64.446 & 42.462 & 62.005 & 1.465 & 26/2 \\
late\_finite\_64 & 328.629 & 184.098 & 327.261 & 1.785 & 66/2 \\
late\_compressed & 127.182 & 112.913 & 128.522 & 1.161 & 25/2 \\
early\_compressed & 127.902 & 117.397 & 127.392 & 1.071 & 25/2 \\
positive\_forest & 115.400 & 115.676 & 112.968 & 1.025 & 32/32 \\
all\_singleton & 36.872 & 12.712 & 12.484 & 2.902 & 256/256 \\
positive\_mixed & 7.053 & 7.205 & 7.203 & 0.976 & 4/4 \\
single\_wider & 1.171 & 1.184 & 1.167 & 1.006 & 2/2 \\
direct\_three & 4.020 & 4.026 & 4.024 & 1.000 & 0/0 \\
\bottomrule
\end{tabular}
\end{center}


The ratio is the median paired old/current elapsed time. All source
grammars, samples, factor orders, proof sizes, projection calls and shared
resource counts are in \path{../fast/results/lazy_forest_20261008.json}.
The experiment includes positive forests, a negative factor early in
strand order, a single wider factor and a direct three-braid control;
the selected examples are not restricted to favorable early exits.


The first tested estimate counted every nonempty node hull. A complete-call
benchmark exposed a scheduling defect: concatenations between disjoint
factors inflated the estimate for early factors, even though they alias
away under projection. On the early-compressed negative fixture, the
producer visited 22 factors instead of the old policy's one, with paired
old/new ratio 0.849. The first measurements are preserved in
\path{../fast/results/lazy_forest_initial_20261008.json}.

The final bound counts possible rule creations, using child-hull
intersections as proved above. It is exact on the tested block-separated
forest grammars, regardless of negative-factor placement. A dedicated
regression checks that both early and late compressed negative factors are
visited first. Nonconvex supports still give a conservative bound, so no
claim of optimal scheduling is made for arbitrary input DAGs.

The final run has 250 measured calls and 50 excluded warmups, taking
24.451 seconds. For 64 positive three-braid factors followed by a finite
negative factor, complete recognition falls from 328.629 to 184.098 ms,
with median paired ratio 1.785. Production plus replay makes two projection
calls instead of 66, and charged work falls from 414,154 to 98,011.
The eight- and 24-positive-factor versions have paired ratios 1.255 and
1.465. The current eager arm remains close to the old implementation on
these cases, isolating preparation as the mechanism.

For the 24-factor compressed-negative fixtures, both early and late
placements now visit exactly one factor and make two projection calls
instead of 25. Charged work is 81,033 instead of 124,353 in both placements.
The paired time ratios are 1.071 for early placement and 1.161 for late
placement. These smaller elapsed-time differences should be read with the
controls: identical-old ratios across the run range from 0.935 to 1.031,
and identical-current ratios from 0.966 to 1.060. The structural correction
is exact even where timing evidence for a small gain is noisy.

The 128-strand chain consisting entirely of singleton generators improves
from 36.872 to 12.712 ms (paired ratio 2.902). It still needs all 128 factors
and 256 producer/replay projection calls. Here the improvement comes from
the canonical one-strand shortcut, not early stopping; current eager
preparation also benefits, with median 12.484 ms. Its charged work falls
from 77,209 to 13,205.

The sixteen-factor positive compressed forest needs every projection.
Its old/new medians are 115.400/115.676 ms, with no systematic gain asserted.
The smaller mixed positive case has paired ratio 0.976; planning costs can
outweigh savings when all factors are needed. Single wider-factor and direct
three-braid controls remain close to unchanged. The implementation therefore
retains an eager policy for comparison and makes no universal claim about
the best discovery order.

\paragraph{Remaining complexity work.}
The complete recognizer now avoids unused factor projections. Its input is
still a supplied grammar. It neither constructs such a
grammar of controlled size from every knot diagram nor bounds the exceptional
parameter for every input. Section~\ref{sec:compact-permutations} subsequently reduces permutation
storage and evaluation work on local forest nodes while preserving the full
component check. Reconciling incoming overlap-search ideas with the maintained
group backend remains a separate next step toward the general objective.
